\section{Cognitive Manifold Alignment of NDH}
\label{sec:ndh-cognitive-alignment}

NDH interacts with cognitive manifolds through alignment rules that synchronize
spectral emission, modal behavior, and recursion structure with an external
cognitive geometry. Cognitive alignment defines NDH's role as a geometric
augmentation of cognitive processes.

\subsection{Cognitive Manifold}

Let $X$ be a cognitive manifold equipped with a cognitive metric $g_X$ and a
cognitive state map
\[
\chi : X \to \Gamma,
\]
where $\Gamma$ is a cognitive state space. The manifold $X$ contains structural
constraints governing cognitive continuity and stability.

\subsection{Alignment Map}

NDH aligns with the cognitive manifold through a smooth map
\[
\Theta : M \to X.
\]
The alignment map preserves modal structure:
\[
\Theta(M_i) \subseteq X_i,
\]
where $X_i$ are modal-aligned cognitive regions.

\subsection{Spectral-Cognitive Synchronization}

Spectral emission synchronizes with cognitive state through:
\[
\chi(\Theta(x)) = G(\{\phi_k(x)\}_{k \in K}),
\]
for some smooth function $G$. This defines how NDH's spectral components map to
cognitive states.

\subsection{Cognitive Continuity}

Cognitive continuity requires:
\[
d_X(\Theta(x), \Theta(y))
  \le L_X \cdot \|\mathcal{N}_{DH}(x) - \mathcal{N}_{DH}(y)\|,
\]
for some constant $L_X$. This ensures that NDH's spectral variation does not
violate cognitive continuity.

\subsection{Cognitive Stability}

Cognitive stability requires that recursion behavior preserves cognitive state:
\[
\chi(\Theta(R_i(x))) = \chi(\Theta(x)).
\]
This defines stability-preserving recursion across modal regions.

\subsection{Alignment of Modal Boundaries}

Modal boundaries map to cognitive boundaries:
\[
\Theta(\mathcal{B}_{ij}) \subseteq \partial X_{ij}.
\]
This ensures that NDH's modal transitions correspond to transitions in cognitive
geometry.

\subsection{Prosthetic Cognitive Effect}

NDH induces a prosthetic cognitive effect defined by:
\[
\text{Aug}(x) = (\mathcal{N}_{DH}(x), \Theta(x)).
\]
The augmentation map $\text{Aug}$ defines NDH's role as a geometric cognitive
signal augmenting the cognitive manifold.

\subsection{Alignment Summary}

Cognitive manifold alignment consists of a smooth alignment map, spectral-cognitive
synchronization, continuity and stability constraints, modal boundary alignment,
and a prosthetic cognitive augmentation. This defines NDH's interaction with
cognitive geometry.
